\section{Endpoint branching at \texorpdfstring{\(y=0\)}{y=0}}\label{app:endpoint-branching}

This appendix is part of the proof of Theorem~\ref{thm:endpoint-zero-quantization}: Proposition~\ref{prop:endpoint-expansion} supplies the branch expansions, amplitudes, and remainder estimate used in Proposition~\ref{prop:endpoint-cosine-hypotheses} and in the fixed-window isolation Proposition~\ref{prop:fixed-window-endpoint-isolation}.
It verifies the hypotheses of Lemma~\ref{lem:two-branch-oscillation-law} for the quartic transfer family at the square-root endpoint $y_c=0$ and records the branch expansion predicted by Lemma~\ref{lem:generic-ep2-branching}.
The coefficient comparisons below give the exact constants \(A=1/\sqrt2\), \(B=5/8\), \(A_*=1/2\), \(B_*=3\sqrt2/16\), and, for each compact \(K\subset\CC\), the estimate \(R_m(-u^2/m^2)=O_K(m^{-2})\).
Fix the disc
\[
D_r:=\{y\in\CC:\abs{y}<r\},\qquad r:=\tfrac14,
\]
which is small enough that $(\lambda,y)=(1,0)$ is the only repeated point of $\Pi(\lambda,y)$ in $D_r$.
Indeed, the other finite branch values are $1$ and the three roots of $c(y)=256y^3+411y^2+165y+32$ by Corollary~\ref{cor:disc-factor}.
Let $\eta,\bar\eta$ be the non-real roots of $c(y)$.
Since the product of the three roots is $-32/256=-1/8$,
\[
y_{\mathrm{sub}}\eta\bar\eta=-\frac18.
\]
Hence
\[
|\eta|^2=\eta\bar\eta=-\frac{1}{8y_{\mathrm{sub}}}
=\frac{1}{8|y_{\mathrm{sub}}|}.
\]
Lemma~\ref{lem:cubic-branch-real-root} gives $-2<y_{\mathrm{sub}}<-1$, so $|y_{\mathrm{sub}}|<2$, and therefore
\[
|\eta|^2>\frac1{16},
\qquad
|\eta|>\frac14.
\]
Thus the two non-real branch values lie outside $D_{1/4}$.
The branch value $1$ and the real branch value $y_{\mathrm{sub}}$ also lie outside $D_{1/4}$, so $0$ is the only branch value in the disc.
All holomorphic expansions below are carried out on $D_r$ (or on $|s|<\sqrt{r}$ when $y=s^2$), and every remainder bound is uniform on that set.
The same radius $r=1/4$, the same branch coordinate $y=s^2$, and the same constants in \eqref{eq:appendix-cosine-constants} are the ones used in Proposition~\ref{prop:endpoint-cosine-hypotheses}; no independent local constants are introduced in Section~\ref{sec:positive-maximal-modulus-branch}.
The body uses Proposition~\ref{prop:endpoint-expansion} in this fixed form: the branch-coordinate domain, the constants \((a,b,A_*,B_*)\), the symmetries \(\lambda_-(s)=\lambda_+(-s)\) and \(C_-(s)=C_+(-s)\), nonvanishing after shrinking inside \(|s|<1/2\), and the estimate \(R_m(-u^2/m^2)=O_K(m^{-2})\) for the compact set \(K\) chosen in the endpoint scale.
The local expansion produces the constants
\begin{equation}\label{eq:appendix-cosine-constants}
\lambda_c=1,\qquad
a=\frac{1}{\sqrt2},\qquad
b=\frac38,\qquad
A_*=\frac12,\qquad
B_*=\frac{3\sqrt2}{16},
\end{equation}
where \(a=A\) is the square-root branch coefficient, \(b=B-A^2/2\) is the logarithmic quadratic coefficient coming from \(B=5/8\), and these constants yield the model-specific scaling law stated in Theorem~\ref{thm:endpoint-zero-quantization}.

\begin{lemma}[Uniform endpoint coefficient computation]\label{lem:endpoint-coefficient-computation}
On the branch-coordinate disc \(|s|<1/2\), the two branches over the double root \(\lambda=1\) are obtained by the convergent ansatz
\[
\lambda_\pm(s)=1\pm As+Bs^2\pm Cs^3+O(s^4),
\qquad y=s^2,
\]
where
\[
A=\frac1{\sqrt2},
\qquad
B=\frac58,
\qquad
C=-\frac{43}{64\sqrt2}
\]
for the plus branch, and the minus branch is obtained by replacing \(s\) by \(-s\).
The two non-colliding branches and their amplitudes satisfy, uniformly for \(|y|<1/4\),
\begin{gather*}
\lambda_{-1}(y)=-1-\frac y4+O(y^2),
\qquad
\lambda_0(y)=-y+O(y^2),\\
C_{-1}(y)=\frac y8+O(y^2),
\qquad
C_{\mathrm{small}}(y)=-y+O(y^2).
\end{gather*}
Finally, the branching amplitudes have the uniquely determined expansion
\[
C_\pm(s)=\frac12\mp\frac{3\sqrt2}{16}s+O(s^2).
\]
All three displayed identities are algebraic consequences of \(\Pi\), \(N\), and \(D\), with constants independent of any compact set later used in the endpoint scaling variable \(u\).
\end{lemma}

\begin{proof}
The preceding paragraph shows that \(|s|<1/2\) maps into the disc \(|y|<1/4\) containing no branch value except \(0\).
Thus the two colliding roots are single-valued holomorphic functions of \(s\), while the two roots near \(-1\) and \(0\) are holomorphic functions of \(y\) on the same disc by the implicit-function theorem.

Write \(w=\lambda-1\).  A direct expansion gives
\begin{equation}\label{eq:endpoint-recursive-polynomial}
\Pi(1+w,s^2)=2w^2+3w^3+w^4-s^2(1+4w+2w^2)+s^4.
\end{equation}
Substitute \(w=As+Bs^2+Cs^3+Ds^4+O(s^5)\).  The coefficients of \(s^2,s^3,s^4\) in \eqref{eq:endpoint-recursive-polynomial} are
\begin{align*}
&2A^2-1,\\
&A(3A^2+4B-4),\\
&A^4+9A^2B-2A^2+4AC+2B^2-4B+1.
\end{align*}
The first equation gives \(A^2=1/2\); choosing \(A=1/\sqrt2\) labels the plus branch.
The second gives \(B=1-3A^2/4=5/8\).
With these values, the third gives
\[
C=\frac{37A^4-56A^2+8}{32A}=-\frac{43}{64\sqrt2}.
\]
Replacing \(s\) by \(-s\) gives the conjugate Puiseux branch and proves the first assertion.

For the non-colliding branches, substitute \(
\lambda=-1+ay+O(y^2)
\) into \(\Pi(\lambda,y)=0\).  The coefficient of \(y\) is \(-4a-1\), hence \(a=-1/4\).
Similarly, substituting \(\lambda=ay+O(y^2)\) gives \(a=-1\).
The residue formula
\[
C_j(y)=\frac{\lambda_j(y)^3N(1/\lambda_j(y),y)}{\partial_\lambda\Pi(\lambda_j(y),y)}
\]
then gives
\[
C_{-1}(y)=\frac{y}{8}+O(y^2),
\qquad
C_{\mathrm{small}}(y)=-y+O(y^2),
\]
after inserting
\[
N(t,y)=1+yt+(y^2-y-1)t^2+(y^3-2y)t^3.
\]
The divisions are legitimate uniformly on \(|y|<1/4\), after reducing the branch labels if necessary, because \(\partial_\lambda\Pi(-1,0)=-4\) and \(\partial_\lambda\Pi(0,0)=1\).

It remains to derive the first amplitude coefficient for the colliding pair.
This is a direct residue computation, not a consequence of the later spectral coefficient comparison.
For \(\epsilon=\pm1\), put
\[
\lambda_\epsilon(s)=1+\epsilon\frac{s}{\sqrt2}+\frac58s^2-
\epsilon\frac{43}{64\sqrt2}s^3+O(s^4),
\qquad y=s^2,
\]
where \(\lambda_{+1}=\lambda_+\) and \(\lambda_{-1}=\lambda_-\).
Substituting this expansion into the numerator and denominator of the residue formula gives
\begin{align}
\lambda_\epsilon(s)^3 N(1/\lambda_\epsilon(s),s^2)
&=
\epsilon\sqrt2\,s+\frac34s^2+
\epsilon\frac{125\sqrt2}{64}s^3+O(s^4),\notag\\
\partial_\lambda\Pi(\lambda_\epsilon(s),s^2)
&=
2\epsilon\sqrt2\,s+3s^2+
\epsilon\frac{105\sqrt2}{32}s^3+O(s^4).
\label{eq:direct-colliding-residue-cancellation-lemma}
\end{align}
Both the numerator and the denominator vanish to first order in the branch coordinate, and the displayed linear terms are nonzero for \(\epsilon=\pm1\).
Dividing the two Taylor series in \eqref{eq:direct-colliding-residue-cancellation-lemma} proves that the apparent pole is removable and gives the holomorphic expansion
\begin{equation}\label{eq:direct-colliding-residue-expansion}
C_\epsilon(s)
=
\frac{\lambda_\epsilon(s)^3N(1/\lambda_\epsilon(s),s^2)}
{\partial_\lambda\Pi(\lambda_\epsilon(s),s^2)}
=
\frac12-
\epsilon\frac{3\sqrt2}{16}s+\frac{7}{16}s^2+O(s^3).
\end{equation}
In particular,
\[
C_+(s)=\frac12-\frac{3\sqrt2}{16}s+O(s^2),
\qquad
C_-(s)=\frac12+\frac{3\sqrt2}{16}s+O(s^2),
\]
and the involution \(s\mapsto -s\) gives the compatible symmetry \(C_-(s)=C_+(-s)\).
The rational identity at \(y=0\) is \(Z(t,0)=1/(1-t)\), so this direct expansion also recovers \(C_+(0)=C_-(0)=1/2\).
For comparison with the notation used below, write
\[
C_\pm(s)=\frac12\mp B_*s+E s^2+O(s^3),
\qquad
B_*=\frac{3\sqrt2}{16},
\qquad
E=\frac{7}{16}.
\]
Using the branch expansion above,
\[
\lambda_\pm(s)^2=1\pm\sqrt2 s+\frac74s^2+O(s^3),
\qquad
\lambda_\pm(s)^3=1\pm\frac{3}{\sqrt2}s+\frac{27}{8}s^2+O(s^3).
\]
The exact coefficients from \(D(t,y)Z(t,y)=N(t,y)\) are \(Z_2(s^2)=1+2s^2+O(s^4)\) and \(Z_3(s^2)=1+3s^2+O(s^4)\).
The non-colliding pair contributes \((1/8)s^2+O(s^4)\) to \(Z_2\) and \(-(1/8)s^2+O(s^4)\) to \(Z_3\).
The following coefficient comparison is therefore only a consistency check for the exact spectral identities.  Comparing the coefficients of \(s^2\) gives
\[
-2\sqrt2 B_*+2E=\frac18,
\qquad
-3\sqrt2 B_*+2E=-\frac14,
\]
which is satisfied by \(B_*=3\sqrt2/16\) and \(E=7/16\).
This proves the amplitude assertion and all bounds are uniform because each expansion is the Taylor expansion of a holomorphic function on the fixed disc.
\end{proof}

\subsection{Partial fractions, local spectral branches, and amplitude derivation}

For $y\in D_r\setminus\{0\}$ the quartic polynomial $\Pi(\lambda,y)$ has four simple roots $\lambda_1(y),\dots,\lambda_4(y)$, and the rational generating function admits the partial-fraction decomposition
\begin{equation}\label{eq:partial-fractions}
Z(t,y)=\sum_{j=1}^4 \frac{C_j(y)}{1-\lambda_j(y)t},
\qquad
C_j(y)=\frac{\lambda_j(y)^3\,N(1/\lambda_j(y),y)}{\partial_\lambda\Pi(\lambda_j(y),y)}.
\end{equation}
Consequently,
\begin{equation}\label{eq:spectral-decomposition}
Z_m(y)=\sum_{j=1}^4 C_j(y)\lambda_j(y)^m.
\end{equation}

\subsubsection*{Puiseux expansion of \texorpdfstring{$\lambda_\pm$}{lambda\_pm} at \texorpdfstring{$y=0$}{y=0}}

Near $y=0$, write $y=s^2$.
The relevant partial derivatives are
\[
\Pi(1,0)=0,\qquad
\Pi_\lambda(1,0)=0,\qquad
\Pi_{\lambda\lambda}(1,0)=4,\qquad
\Pi_y(1,0)=-1.
\]
Lemma~\ref{lem:generic-ep2-branching} with $P=\Pi$, $\lambda_c=1$, $y_c=0$ gives $\alpha^2=-2P_y/P_{\lambda\lambda}\big|_{(1,0)}=1/2$, so $\alpha=1/\sqrt2$.
Hence the two colliding branches have leading behavior
\[
\lambda_\pm(s)=1\pm \frac{s}{\sqrt2}+O(s^2).
\]
To obtain the next coefficients, write
\[
\lambda_\pm(s)=1\pm As+Bs^2\pm Cs^3+O(s^4),
\qquad \lambda_-(s)=\lambda_+(-s),
\]
where the symmetry is guaranteed by Lemma~\ref{lem:generic-ep2-branching}.
The coefficients $A,B,C$ are determined by expanding
\[
\Pi(1\pm As+Bs^2\pm Cs^3,s^2)=0
\]
and collecting powers of $s$.
Proposition~\ref{prop:endpoint-puiseux} carries out this expansion via the equivalent parametrization $\lambda_{\mathrm{dom}}(y)$ and yields the plus-branch coefficient of $s^3$ equal to $-43/(64\sqrt2)$.
For the reader's convenience, the key step at order $s^3$ is as follows.
Setting $\lambda=1\pm As+Bs^2$ and $y=s^2$ in
\[
\Pi(\lambda,y)
=\lambda^4-\lambda^3-(2y+1)\lambda^2+\lambda+y(y+1)
\]
gives the following coefficient comparisons:
the $s^2$ coefficient gives $2A^2-1=0$ (confirming $A=1/\sqrt2$);
the $s^3$ coefficient (for the $\pm$ branch) is $\pm(4AB-5A^3)=0$, giving $B=5A^2/4=5/8$.
The $s^4$ coefficient then determines the plus-branch coefficient $-43/(64\sqrt2)$.
This gives
\begin{equation}\label{eq:lambda-plus-minus}
\lambda_\pm(s)
=
1\pm \frac{1}{\sqrt2}s+\frac58 s^2 \mp \frac{43}{64\sqrt2}s^3+O(s^4).
\end{equation}

\subsubsection*{Logarithmic form}

To verify the exponential form, compute $\log\lambda_\pm(s)$ from \eqref{eq:lambda-plus-minus}:
\[
\log\lambda_\pm(s)
=\log\bigl(1\pm\tfrac{s}{\sqrt2}+\tfrac58s^2\mp\tfrac{43}{64\sqrt2}s^3+O(s^4)\bigr).
\]
Using $\log(1+u)=u-u^2/2+u^3/3+\cdots$ with $u=\pm s/\sqrt2+O(s^2)$:
\[
\log\lambda_\pm(s)=\pm\frac{s}{\sqrt2}+\frac38 s^2+O(s^3),
\]
which matches \eqref{eq:log-lambda-plus-minus} and confirms
\begin{equation}\label{eq:lambda-pm-exp}
\lambda_\pm(s)=\exp\!\Bigl(\pm\frac{s}{\sqrt2}+\frac38 s^2+O(s^3)\Bigr).
\end{equation}

\subsubsection*{Explicit derivation of the amplitude coefficients \texorpdfstring{$C_\pm$}{C\_pm}}

The simple roots at $\lambda=-1$ and $\lambda=0$ give analytic branches $\lambda_{-1}(y)$ and $\lambda_0(y)$ by the implicit function theorem, since $\partial_\lambda\Pi(-1,0)=-4\neq0$ and $\partial_\lambda\Pi(0,0)=1\neq0$.
Shrinking $r$ if necessary, direct Taylor expansion of $\Pi(\lambda,y)=0$ near each gives
\begin{equation}\label{eq:other-branches}
\lambda_{-1}(y)=-1-\frac{y}{4}+O(y^2),
\qquad
\lambda_0(y)=-y+O(y^2),
\end{equation}
uniformly for $|y|<r_0$.
The residue formula \eqref{eq:partial-fractions} is algebraic in $(\lambda_j,y)$, so the amplitudes $C_{-1}(y)$ and $C_{\mathrm{small}}(y)$ are holomorphic on the same disc:
\begin{equation}\label{eq:other-amplitudes}
C_{-1}(y)=\frac{y}{8}+O(y^2),
\qquad
C_{\mathrm{small}}(y)=-y+O(y^2),
\end{equation}
uniformly for $|y|<r_0$.
All $O$-bounds below hold in this complex disc.

\begin{lemma}[Plus-branch residue cancellation]\label{lem:plus-branch-residue-cancellation}
With $y=s^2$ and the plus branch
\[
\lambda_+(s)=1+\frac{s}{\sqrt2}+\frac58s^2-\frac{43}{64\sqrt2}s^3+O(s^4),
\]
the residue
\[
C_+(s)=
\frac{\lambda_+(s)^3N(1/\lambda_+(s),s^2)}
{\partial_\lambda\Pi(\lambda_+(s),s^2)}
\]
is holomorphic at $s=0$ and satisfies
\[
C_+(s)=\frac12-\frac{3\sqrt2}{16}s+O(s^2).
\]
\end{lemma}

\begin{proof}
For $s\ne0$, the two colliding roots are simple and the displayed formula is the usual residue coefficient from \eqref{eq:partial-fractions}.
Direct expansion gives
\begin{align}
\lambda_+(s)^3 N(1/\lambda_+(s),s^2)
&=
\sqrt2\,s+\frac34s^2+O(s^3),\notag\\
\partial_\lambda\Pi(\lambda_+(s),s^2)
&=
2\sqrt2\,s+3s^2+O(s^3).
\label{eq:plus-residue-expansion-cancellation}
\end{align}
Both numerator and denominator therefore vanish to first order in $s$, and the first coefficient in the denominator is nonzero.
After cancelling the common factor $s$, the quotient extends holomorphically through $s=0$.
Dividing the two Taylor series in \eqref{eq:plus-residue-expansion-cancellation} gives
\[
C_+(s)
=
\frac{\sqrt2+\frac34s+O(s^2)}{2\sqrt2+3s+O(s^2)}
=
\frac12-\frac{3\sqrt2}{16}s+O(s^2),
\]
as claimed.
\end{proof}

The amplitudes $C_\pm(s)$ for the branching pair are holomorphic in $s$ away from $s=0$ because the two roots are simple there, so \eqref{eq:partial-fractions} gives the unique residue coefficients of the partial-fraction expansion.
Lemma~\ref{lem:plus-branch-residue-cancellation} proves the removability and first coefficient for the plus branch.
Applying the same calculation to the branch \(\lambda_-(s)=\lambda_+(-s)\), or equivalently replacing $s$ by $-s$, gives
\begin{equation}\label{eq:colliding-residue-laurent-cancellation}
C_\pm(s)
=
\frac12\mp\frac{3\sqrt2}{16}s+O(s^2).
\end{equation}
Thus the possible Laurent pole has zero coefficient; the individual residues, not merely their symmetric sum, extend holomorphically through $s=0$.
The same calculation also proves the symmetry-compatible identities $C_-(s)=C_+(-s)$ and $C_\pm(0)=1/2$.

We now record how the residue-derived constants fit the exact coefficient identities.

\textbf{Leading coefficient $A_*$.}
At $y=0$: direct factorization gives $D(t,0)=1-t-t^2+t^3=(1-t)^2(1+t)$ and $N(t,0)=1-t^2=(1-t)(1+t)$, so $Z(t,0)=1/(1-t)$ and $Z_m(0)=1$ for all $m\ge0$.
The direct residue expansion \eqref{eq:colliding-residue-laurent-cancellation} gives the individual removable limits $C_+(0)=C_-(0)=1/2$; equivalently, from the spectral identity at $s=0$, using $\lambda_\pm(0)=1$ and $C_{-1}(0)=C_{\mathrm{small}}(0)=0$:
\[
1 = Z_m(0) = 2C_+(0) \implies C_+(0) = C_-(0) = \tfrac{1}{2},\quad A_*=\tfrac12.
\]

\textbf{First-order coefficient $B_*$.}
The direct expansion \eqref{eq:colliding-residue-laurent-cancellation} gives
\[
C_\pm(s)=\tfrac12\mp B_*s+O(s^2),
\qquad
B_*=\frac{3\sqrt2}{16}.
\]
Write more finely as $C_\pm(s)=\tfrac12\mp B_*s+E\,s^2+O(s^3)$ for some coefficient $E$.
The order-$s$ terms in the spectral sum $C_+\lambda_+^m+C_-\lambda_-^m$ cancel identically for every $m$, so $B_*$ first appears at order~$s^2$.
The comparison of the $s^2$ coefficients for the two initial values $m=2$ and $m=3$ is only a consistency check against the exact spectral identities.
These two values are enough to recover the same $B_*$ independently because, after $A_*=1/2$ is fixed, the order-$s^2$ truncation has exactly the two unknowns $B_*$ and $E$.
The recurrence gives $Z_2(s^2)=(1+s^2)^2=1+2s^2+O(s^4)$ and $Z_3(s^2)=(1+s^2)^3=1+3s^2+O(s^4)$.

For the two colliding branches, the required order-$s^2$ contributions are:
\begin{align*}
\lambda_\pm(s)^2
&=1\pm\sqrt2\,s+\frac74s^2+O(s^3),\\
C_+(s)\lambda_+(s)^2+C_-(s)\lambda_-(s)^2
&=1+\left(\frac74-2\sqrt2\,B_*+2E\right)s^2+O(s^3),\\
Z_2(s^2)
&=1+2s^2+O(s^4),
\end{align*}
and
\begin{align*}
\lambda_\pm(s)^3
&=1\pm\frac{3}{\sqrt2}s+\frac{27}{8}s^2+O(s^3),\\
C_+(s)\lambda_+(s)^3+C_-(s)\lambda_-(s)^3
&=1+\left(\frac{27}{8}-3\sqrt2\,B_*+2E\right)s^2+O(s^3),\\
Z_3(s^2)
&=1+3s^2+O(s^4).
\end{align*}
Here the non-colliding entries come directly from \eqref{eq:other-branches}--\eqref{eq:other-amplitudes}:
\[
C_{-1}(s^2)\lambda_{-1}(s^2)^2=\frac18s^2+O(s^4),
\qquad
C_{\mathrm{small}}(s^2)\lambda_0(s^2)^2=O(s^6),
\]
and
\[
C_{-1}(s^2)\lambda_{-1}(s^2)^3=-\frac18s^2+O(s^4),
\qquad
C_{\mathrm{small}}(s^2)\lambda_0(s^2)^3=O(s^8).
\]
Equating the $s^2$ coefficient in the $m=2$ column would require
\begin{equation}\label{eq:Bstar-m2}
\frac74 - 2\sqrt2\,B_* + 2E + \frac18 = 2,
\qquad\text{i.e.,}\qquad
-2\sqrt2\,B_* + 2E = \frac18.
\end{equation}
Equating the $s^2$ coefficient in the $m=3$ column would require
\begin{equation}\label{eq:Bstar-m3}
\frac{27}{8} - 3\sqrt2\,B_* + 2E - \frac18 = 3,
\qquad\text{i.e.,}\qquad
-3\sqrt2\,B_* + 2E = -\frac14.
\end{equation}
The direct residue expansion \eqref{eq:direct-colliding-residue-expansion} gives $B_*=3\sqrt2/16$ and $E=7/16$; substituting these two values into \eqref{eq:Bstar-m2} and \eqref{eq:Bstar-m3} verifies both identities.
For later reference we retain the first-order consequence
\begin{equation}\label{eq:C-plus-minus}
C_\pm(s)
=
\frac12\mp \frac{3\sqrt2}{16}\,s+O(s^2).
\end{equation}

\textbf{Cross-check.}
The same algebra for $m=1$ (where $Z_1(s^2)=1+s^2$) gives another check: the maximal-modulus-branch sum contributes $\tfrac58-\sqrt2\,B_*+2E$ to the $s^2$ coefficient, the $\lambda_{-1}$ branch contributes $-\tfrac18\,s^2$, and the total must equal~$1$.
With the residue-derived values $B_*=3\sqrt2/16$ and $E=7/16$, this identity is also satisfied.

The two branches issuing from $\lambda_c=1$ therefore have the exponential form \eqref{eq:lambda-pm-exp} required by Lemma~\ref{lem:two-branch-oscillation-law}, with constants $A_*=1/2$ and $B_*=3\sqrt2/16$.

\subsection{Endpoint-scale remainder bound}

Fix a compact set \(K\subset\CC\) and let \(M_K:=\sup_{u\in K}\abs{u}\).
Set
\[
y=-\frac{u^2}{m^2},
\qquad u\in K.
\]
Then \(\abs{y}\le M_K^2m^{-2}\), so for all sufficiently large \(m\) the point \(y\) lies in the complex disc \(|y|<r_0\).
On this disc the holomorphic expansions \eqref{eq:other-branches}--\eqref{eq:other-amplitudes} are valid.
In particular, there exists a constant \(C_K>0\) (depending only on \(K\)) such that, for all \(u\in K\) and all sufficiently large \(m\),
\begin{equation}\label{eq:y-bound}
\abs{\lambda_{-1}(y)+1}\le \frac{C_K}{m^2},
\quad
\abs{C_{-1}(y)}\le \frac{C_K}{m^2},
\quad
\abs{\lambda_0(y)}\le \frac{C_K}{m^2},
\quad
\abs{C_{\mathrm{small}}(y)}\le \frac{C_K}{m^2}.
\end{equation}

Define \(h(y):=-\lambda_{-1}(y)\).
Since \(h(0)=1\) and \(h(y)=1+y/4+O(y^2)\) holds uniformly on \(|y|<r_0\), the function \(h\) is nonvanishing for \(|y|\) sufficiently small, and a holomorphic branch of the logarithm is well defined:
\[
\log h(y)=\frac{y}{4}+O(y^2),
\qquad |y|<r_0.
\]
Hence, for the holomorphic branch of \(\log\lambda_{-1}\) near \(y=0\),
\begin{equation}\label{eq:log-lambda-minus-one}
\log \lambda_{-1}(y)=i\pi+\frac{y}{4}+O(y^2),
\qquad |y|<r_0.
\end{equation}
Exponentiating and substituting \(y=-u^2/m^2\) gives
\[
\lambda_{-1}\!\left(-\tfrac{u^2}{m^2}\right)^m
=
(-1)^m\exp\!\left(-\frac{u^2}{4m}+O_K(m^{-3})\right).
\]
In particular, \(\bigl|\lambda_{-1}(-u^2/m^2)^m\bigr|=\exp(O_K(m^{-1}))=1+O_K(m^{-1})\), and therefore
\begin{equation}\label{eq:minus-one-contribution-bound}
\left|
C_{-1}\!\left(-\frac{u^2}{m^2}\right)
\lambda_{-1}\!\left(-\frac{u^2}{m^2}\right)^m
\right|
\le C_K m^{-2}e^{C_K/m}
=
O_K(m^{-2}),
\end{equation}
uniformly for all $u\in K$.

For the root near \(0\), the holomorphic expansions \eqref{eq:other-branches}--\eqref{eq:other-amplitudes} give
\[
\lambda_0(y)=-y+O(y^2),
\qquad
C_{\mathrm{small}}(y)=-y+O(y^2),
\qquad |y|<r.
\]
Substituting \(y=-u^2/m^2\) and using $|\lambda_0(-u^2/m^2)|\le C_K m^{-2}$:
\begin{equation}\label{eq:small-root-contribution-bound}
\left|
C_{\mathrm{small}}\!\left(-\frac{u^2}{m^2}\right)
\lambda_0\!\left(-\frac{u^2}{m^2}\right)^m
\right|
\le
C_K m^{-2}\bigl(C_K m^{-2}\bigr)^m
=
O_K(m^{-2m-2}),
\end{equation}
uniformly for all $u\in K$.

Writing
\[
R_m(y):=
C_{-1}(y)\lambda_{-1}(y)^m+C_{\mathrm{small}}(y)\lambda_0(y)^m,
\]
we conclude that
\begin{equation}\label{eq:appendix-remainder-bound}
R_m\!\left(-\frac{u^2}{m^2}\right)=O_K(m^{-2})
\end{equation}
uniformly for \(u\in K\).
Equivalently, for each compact \(K\) there is a constant \(M_K\) such that
\[
\abs{R_m(-u^2/m^2)}\le M_K\lambda_c^m m^{-2}
\qquad (u\in K),
\]
because the endpoint constant in \eqref{eq:appendix-cosine-constants} is \(\lambda_c=1\).
This is precisely the remainder estimate required by Lemma~\ref{lem:two-branch-oscillation-law}; Proposition~\ref{prop:fixed-window-endpoint-isolation} specifies the compact sets on which it is applied.

\subsection{Specialization of the local normal form}

Fix again a compact set $K\subset\CC$.
The logarithm of \eqref{eq:lambda-plus-minus} has the expansion
\begin{equation}\label{eq:log-lambda-plus-minus}
\log \lambda_\pm(s)
=
\pm \frac{1}{\sqrt2}s+\frac38 s^2+O(s^3).
\end{equation}
With $s=iu/m$, one obtains
\begin{equation}\label{eq:lambda-power}
\lambda_\pm\!\left(\frac{iu}{m}\right)^m
=
\exp\!\left(
\pm \frac{iu}{\sqrt2}-\frac{3u^2}{8m}+O_K(m^{-2})
\right).
\end{equation}
Similarly, \eqref{eq:C-plus-minus} becomes
\begin{equation}\label{eq:amplitude-scale}
C_\pm\!\left(\frac{iu}{m}\right)
=
\frac12\mp i\frac{3\sqrt2}{16}\frac{u}{m}+O_K(m^{-2}).
\end{equation}
Therefore the constants in Lemma~\ref{lem:two-branch-oscillation-law} are
\begin{equation}\label{eq:appendix-cosine-constants-repeated}
a=\frac{1}{\sqrt2},
\qquad
b=\frac38,
\qquad
A_*=\frac12,
\qquad
B_*=\frac{3\sqrt2}{16}.
\end{equation}
Substituting \eqref{eq:lambda-power} and \eqref{eq:amplitude-scale} into \eqref{eq:spectral-decomposition} gives
\begin{align}
Z_m\!\left(-\frac{u^2}{m^2}\right)
&=
\exp\!\left(-\frac{3u^2}{8m}\right)
\Biggl[
\left(\frac12-i\frac{3\sqrt2}{16}\frac{u}{m}\right)e^{iu/\sqrt2}
\notag\\
&\qquad\qquad
+\left(\frac12+i\frac{3\sqrt2}{16}\frac{u}{m}\right)e^{-iu/\sqrt2}
+O_K(m^{-2})
\Biggr]
\notag\\
&=
\exp\!\left(-\frac{3u^2}{8m}\right)
\left[
\cos\!\left(\frac{u}{\sqrt2}\right)
+\frac{3u}{4\sqrt2\,m}\sin\!\left(\frac{u}{\sqrt2}\right)
+O_K(m^{-2})
\right].\label{eq:appendix-scaling}
\end{align}
Using \(\cos(\theta-\delta)=\cos\theta+\delta\sin\theta+O(\delta^2)\) with
\(\delta=3u/(4\sqrt2\,m)\),
one obtains the uniform scaling law
\begin{equation}\label{eq:appendix-shifted-cosine}
Z_m\!\left(-\frac{u^2}{m^2}\right)
=
\exp\!\left(-\frac{3u^2}{8m}\right)
\left[
\cos\!\left(\frac{u}{\sqrt2}-\frac{3u}{4\sqrt2\,m}\right)
+ O_K(m^{-2})
\right].
\end{equation}
This identity records the model cosine scale supplied by the local constants.

\begin{proposition}[Endpoint local expansion proposition]\label{prop:endpoint-expansion}
For the quartic transfer family, after writing \(y=s^2\) on the fixed disc \(D_{1/4}\), the two branches issuing from the double root \(\lambda=1\) and their partial-fraction amplitudes may be labelled so that \(\lambda_-(s)=\lambda_+(-s)\), \(C_-(s)=C_+(-s)\), and
\[
\lambda_\pm(s)
=
1\pm \frac{s}{\sqrt2}+\frac58s^2\mp\frac{43}{64\sqrt2}s^3+O(s^4),
\qquad
C_\pm(s)=\frac12\mp\frac{3\sqrt2}{16}s+O(s^2).
\]
After possibly shrinking only the branch-coordinate disc inside \(|s|<1/2\), these two branches and amplitudes are nonzero.
Moreover, for every compact set \(K\subset\CC\), the sum of the two remaining spectral contributions satisfies
\[
R_m\!\left(-\frac{u^2}{m^2}\right)=O_K(m^{-2})
\]
uniformly for \(u\in K\).
The compact set \(K\) is fixed before the estimate is applied, and the implied constant may depend on it.
\end{proposition}

\begin{proof}
The preceding subsections establish the identities used below; we collect them to make the logical dependence explicit.

For the colliding branches, Lemma~\ref{lem:generic-ep2-branching} applies because
\[
\Pi(1,0)=\Pi_\lambda(1,0)=0,
\qquad
\Pi_{\lambda\lambda}(1,0)=4,
\qquad
\Pi_y(1,0)=-1.
\]
Thus \(y=s^2\) gives two holomorphic branches with leading coefficient \(A^2=1/2\), and the square-root involution gives \(\lambda_-(s)=\lambda_+(-s)\).
Substitution of
\[
\lambda=1\pm As+Bs^2\pm Cs^3+O(s^4)
\]
into \(\Pi(\lambda,s^2)=0\) gives at orders \(s^2\) and \(s^3\)
\[
2A^2-1=0,
\qquad
4AB-5A^3=0,
\]
hence \(A=1/\sqrt2\) and \(B=5/8\); the next coefficient comparison gives the plus-branch coefficient \(-43/(64\sqrt2)\).
This is exactly \eqref{eq:lambda-plus-minus}.

For the amplitudes, the direct residue expansion \eqref{eq:direct-colliding-residue-expansion} proves that the individual residue formulae have removable singularities at $s=0$ and gives holomorphic amplitudes in the same branch coordinate:
\[
C_+(s)=\frac12-\frac{3\sqrt2}{16}s+\frac{7}{16}s^2+O(s^3),
\qquad
C_-(s)=\frac12+\frac{3\sqrt2}{16}s+\frac{7}{16}s^2+O(s^3).
\]
Thus \(C_+(0)=C_-(0)=1/2\), the residue formula is compatible with the same interchange of the two square-root branches, and the labels satisfy \(C_-(s)=C_+(-s)\).
In particular both amplitudes are nonzero after shrinking the branch-coordinate disc.
Writing
\[
C_\pm(s)=\frac12\mp B_*s+E s^2+O(s^3)
\]
therefore has \(B_*=3\sqrt2/16\) and \(E=7/16\).
The exact spectral identities for \(Z_2(s^2)=1+2s^2+O(s^4)\) and \(Z_3(s^2)=1+3s^2+O(s^4)\), including the noncolliding contributions from \eqref{eq:other-branches}--\eqref{eq:other-amplitudes}, then satisfy the two displayed linear equations \eqref{eq:Bstar-m2} and \eqref{eq:Bstar-m3}; this is only the consistency check recorded in \eqref{eq:C-plus-minus}, not an independent derivation of \(B_*\) or of removability.

For the remainder, fix once and for all a compact \(K\subset\CC\) and put \(y=-u^2/m^2\).
The expansions \eqref{eq:other-branches}--\eqref{eq:other-amplitudes} give, uniformly for \(u\in K\),
\begin{gather*}
C_{-1}(y)=O_K(m^{-2}),\\
\lambda_{-1}(y)^m
=(-1)^m\exp\!\left(-\frac{u^2}{4m}+O_K(m^{-3})\right)=O_K(1),
\end{gather*}
and
\[
C_{\mathrm{small}}(y)=O_K(m^{-2}),
\qquad
\lambda_0(y)=O_K(m^{-2}).
\]
Therefore
\begin{gather*}
C_{-1}(y)\lambda_{-1}(y)^m=O_K(m^{-2}),\\
C_{\mathrm{small}}(y)\lambda_0(y)^m=O_K(m^{-2m-2}),
\end{gather*}
which proves \eqref{eq:appendix-remainder-bound}.

Finally, \eqref{eq:lambda-plus-minus} gives
\[
\log\lambda_\pm(s)=\pm\frac{s}{\sqrt2}+\frac38s^2+O(s^3).
\]
With \(s=iu/m\), this is \eqref{eq:lambda-power}; the amplitude expansion becomes \eqref{eq:amplitude-scale}.
Substituting these two uniform expansions and the remainder bound into \eqref{eq:spectral-decomposition} gives the displayed local two-branch estimate \eqref{eq:appendix-scaling}, which is the form passed to Proposition~\ref{prop:endpoint-cosine-hypotheses}.
The subsequent elementary identity
\[
\cos(\theta-\delta)=\cos\theta+\delta\sin\theta+O(\delta^2),
\qquad
\delta=\frac{3u}{4\sqrt2\,m},
\]
is uniform on compact \(K\), and converts \eqref{eq:appendix-scaling} into \eqref{eq:appendix-shifted-cosine}.
\end{proof}

\subsection{Extraction of the first zeros}

Fix $k\ge1$ and define $q_k=(2k-1)\pi/2$.
Applying Lemma~\ref{lem:two-branch-oscillation-law} with
\[
a=\frac{1}{\sqrt2},
\qquad
A_*=\frac12,
\qquad
B_*=\frac{3\sqrt2}{16},
\]
gives a unique zero $u_{m,k}$ of the exact function $u\mapsto Z_m(-u^2/m^2)$ in a small neighborhood of \(\sqrt2\,q_k\), and
\[
u_{m,k}
=
\sqrt2\,q_k\left(1+\frac{3}{4m}\right)+O_k(m^{-2})
=
\frac{\pi(2k-1)}{\sqrt2}\left(1+\frac{3}{4m}\right)+O_k(m^{-2}).
\]
Finally,
\[
y_{m,k}=-\frac{u_{m,k}^2}{m^2}
=
-\frac{\pi^2(2k-1)^2}{2m^2}
-\frac{3\pi^2(2k-1)^2}{4m^3}
+O_k(m^{-4}),
\]
which is the quantization formula stated in Theorem~\ref{thm:endpoint-zero-quantization}.
